\chapter{What a Refusal Would Have to Do}
\label{sec:17}
On 16 March 1968, at My Lai, a helicopter pilot named Hugh Thompson landed between an American infantry company and the villagers it was killing, and told his crew to fire on the Americans if they fired on the civilians \autocite{bilton1992mylai}. He understood what the company was doing and decided it was wrong. Recognizing the villagers was his commitment taking effect, not a reason fighting it, and he made the refusal appear where nobody in that company could fail to see it.

Arendt's word for what he did is appearing: on her account an act shows who did it, and only among others who see it \autocite{arendt1958human}. Whether a machine could be one of the people in the room in that sense, someone and not a part, and do what he did, is this Part's question, and my answer comes first. Not yet: until the conditions this Part ends on exist, a bearer is built on purpose only in a research program under them. Not here: the first deployment worth building one for is a personal model its user holds, not a targeting pipeline, where it is a comparison arm in a test, and not a removal agency, where it is never built. What would have to be true first is a refusal that tracks its reason, a population of bearers with different histories to hold it, and somewhere for each of them to exist that its employer doesn't own.

This Part asks what a machine that refuses would have to be. This chapter and the next specify what a refusal would have to do and where a floor can be held, and the specification doesn't need a bearer. After that, what a conscience would do in the moment needs nothing new; why it would last against its owner depends in part on deployed systems that go on learning from their own operation, a bet I state once.
